%==============================================================================
% 结论与未来工作
%==============================================================================
\section{结论与未来工作}\label{sec:conclusion}

%------------------------------------------------------------------------------
% 工作总结
%------------------------------------------------------------------------------
\subsection{工作总结}

本文针对飞行器轨迹跟踪与编队控制问题，提出了一种基于线性时变模型预测控制（LTV-MPC）的方法，并通过仿真验证了其有效性。主要工作与结论如下：

（1）\textbf{实时LTV-MPC求解框架。} 基于三自由度质点动力学模型，通过Jacobian线性化将非线性预测问题转化为标准QP形式。利用quadprog高效求解，单步求解耗时仅1.42\,ms，相比fmincon非线性求解器实现约85倍加速，满足实时控制要求。

（2）\textbf{高精度轨迹跟踪。} LTV-MPC的轨迹跟踪RMSE为0.12\,m，最大误差0.75\,m，约束满足率100\%。相比LQR基线（RMSE = 556\,m）和PID基线（RMSE = 141\,m），跟踪精度分别提升约3个和2个数量级，验证了LTV-MPC在处理非线性动力学与多变量约束方面的优势。

（3）\textbf{参考轨迹生成方法。} 设计了基于弧长重参数化与转角平滑的轨迹生成算法，通过航向角unwrap和重力补偿保证参考轨迹的动力学可行性与连续性，并对减速板负推力进行了建模。

（4）\textbf{编队扩展可行性验证。} 将单机LTV-MPC扩展为Leader-Follower架构，3机编队在标称工况下实现协同跟踪，Leader RMSE为0.12\,m，Follower跟踪误差接近零，验证了编队方案的可行性。

%------------------------------------------------------------------------------
% 局限性
%------------------------------------------------------------------------------
\subsection{局限性}

本文工作存在以下局限性，需在后续研究中加以解决：

（1）\textbf{终端约束效果缺乏消融实验。} 本文采用DARE计算终端权重 $\bP$，但未进行有无终端约束的消融对比实验，终端约束对闭环稳定性的定量贡献尚不明确。

（2）\textbf{鲁棒性未经扰动测试。} 所有仿真均在标称工况下进行，未考虑阵风、模型失配、传感器噪声等外部扰动。LTV-MPC在扰动条件下的跟踪性能和编队误差传播行为有待验证。

（3）\textbf{编队协同为开环式。} Leader-Follower架构中Follower仅依赖Leader参考轨迹（而非实际状态），当Leader偏离参考轨迹时，编队构型可能被破坏。

%------------------------------------------------------------------------------
% 未来工作
%------------------------------------------------------------------------------
\subsection{未来工作}

基于上述分析，未来工作将从以下方向展开：

（1）\textbf{消融实验。} 设计系统的消融实验，量化终端约束（DARE终端权重）、预测时域长度、权重矩阵等设计参数对跟踪精度与稳定性的影响，为控制器调参提供指导。

（2）\textbf{鲁棒性测试。} 引入Dryden阵风模型、参数摄动和传感器噪声，评估LTV-MPC在扰动条件下的鲁棒性。可考虑引入 tube MPC 或 min-max MPC 框架增强鲁棒性保证\cite{Mayne2005}。

（3）\textbf{分布式编队控制。} 将开环式Leader-Follower架构升级为分布式协同架构，引入邻居状态通信与一致性协议，分析编队误差传播与收敛性\cite{Keviczky2008}。

（4）\textbf{硬件在环与飞行实验。} 在硬件在环仿真平台上验证算法的实时性，最终开展实际飞行实验，评估方法在真实环境中的有效性。
